\section{서론}

요즘 로봇 분야에서는 Physical AI, embodied AI, 로봇 기초모델이라는 말이 자주 함께 등장한다.
산업 담론에서는 Physical AI를 실제 환경에서 지각하고 계획하고 행동하는 기계 지능으로 설명한다 \cite{nvidia_physical_ai_2025}.
로봇 기초모델이라는 표현도 비슷한 기대 속에서 쓰이지만, 이 글은 그런 모델의 성능이나 연구사를 직접 다루지 않는다.
처음에는 카메라와 센서로 보고, 판단하고, 손이나 팔로 실제 물체를 움직이려는 AI 로봇 흐름 정도로 이해하면 충분하다.
다만 이 글의 관심은 AI가 목표를 얼마나 잘 고르는가가 아니라, 그 목표가 실제 접촉으로 이어질 때 힘과 운동을 어떻게 다루는가에 있다.
로봇이 실제 조작 과제로 옮겨갈수록, 모델이 어떤 방식으로 목표를 만들든 마지막 물리 상호작용 단계에서는 접촉 중 힘과 움직임을 다루는 문제가 별도의 학습 고비로 남는다.

이때 로봇은 더 이상 분리된 펜스 안에서 정해진 위치만 반복해서 찍는 장치에 머무르기 어렵다.
사람이 건네는 물체를 잡고, 공구를 표면에 대고, 조립 부품을 끼우고, 문이나 서랍처럼 상태가 변하는 물체를 다루어야 한다.
예를 들어 물류 로봇이 찌그러진 상자를 들어 올린다고 생각해보자.
카메라나 깊이센서는 상자의 찌그러진 모양을 보고, 제어기는 어디를 잡을지 정하며, 손끝은 힘 센서, 촉각 센서, 또는 힘 추정을 이용해 상자를 너무 세게 눌러 망가뜨리지 않으면서도 미끄러지지 않을 만큼 힘을 조절해야 한다.
잡을 위치를 잘 골라도 손끝이 너무 딱딱하게 밀고 들어가면 상자는 찌그러지고, 너무 부드럽게 닿으면 미끄러질 수 있다.
이 장면에서는 인식, 판단, 접촉 힘 조절이 한꺼번에 묶인다.

로봇 매니퓰레이터 제어를 처음 배울 때 가장 익숙한 목표는 원하는 위치로 정확히 이동시키는 것이다.
위치제어는 목표 위치와 현재 위치의 차이를 줄이도록 명령을 만드는 방식이다.
관절각 목표값을 정하고, 현재 관절각과 목표 관절각의 오차를 줄이는 제어기는 직관적이고 구현도 쉽다.
전통적인 자동화 셀을 떠올리면 이 대비가 더 선명하다.
PLC는 공장 장비의 순서와 입출력 신호를 관리하는 산업용 제어 장치이고, 인터록은 특정 안전 조건이 맞지 않으면 다음 동작을 막는 논리이다.
즉 PLC는 셀의 순서와 조건을 관리하는 쪽에 가깝고, 로봇 팔 끝의 접촉력을 직접 세밀하게 조절하는 장치라고 보기는 어렵다.
PLC와 인터록이 상위 순서와 안전 조건을 관리하고, 로봇 컨트롤러가 정해진 위치 궤적을 반복하던 많은 전통적 또는 단순 반복 자동화 셀에서는 이런 방식만으로도 많은 작업을 처리할 수 있었다.
대상이 고정되어 있고, 경로가 반복되며, 비접촉이거나 접촉 상황이 충분히 예측 가능한 작업이 많았기 때문이다.
물론 현대 산업용 로봇 컨트롤러도 힘/토크 센서, compliance 기능, 탐색 동작, 비전 보정, 안전 기능을 제공하는 경우가 있다.
다만 이 대비는 산업 현장이 단순하다는 뜻이 아니다.
말하고 싶은 점은 위치 오차를 줄이는 사고만으로는 접촉 작업을 설명하기 어렵다는 것이다.
그러나 로봇이 딱딱하거나 모델링되지 않은 책상, 벽, 공구, 작업물과 접촉하기 시작하면 위치만을 강하게 추종하는 방식은 위험하거나 불안정해질 수 있다.
사람과 접촉하는 작업에서는 여기에 속도/힘 제한, 안전 표준, 센싱, 비상 정지 같은 별도의 안전 제약도 함께 필요하다.
작은 위치 오차가 큰 접촉력으로 변하고, 이 접촉력이 다시 로봇과 환경의 동역학을 흔들기 때문이다.
숫자로 보면 이 위험은 더 또렷해진다.
스프링처럼 생각하면, 덜 밀리는 물체일수록 강성 $k$가 크다.
준정적 선형 스프링 근사에서 끝단 방향의 등가 강성을 $k_{\mathrm{eff}}$로 두자.
끝단이 벽 안쪽으로 $1~\mathrm{cm}$ 더 들어가려 하고, 그 방향의 등가 강성이 $1000~\mathrm{N/m}$라면 스프링처럼 계산한 힘은
\begin{equation}
  |f|\approx k_{\mathrm{eff}}|\Delta x|=1000\cdot0.01=10~\mathrm{N}
\end{equation}
이 된다.
위치 오차는 겨우 $1~\mathrm{cm}$이지만, 접촉 상황에서는 사람이 손으로 분명히 느낄 만한 힘으로 바뀐다.
이 계산은 동적 충돌, 마찰, 제어 입력 포화, 실제 환경 강성을 모두 제외한 직관용 계산이다.
그래도 작은 위치 오차가 접촉력으로 바뀐다는 사실은 충분히 보여준다.
힘 방향은 작용과 반작용으로 구분하면 된다.
로봇이 오른쪽 벽을 오른쪽으로 밀면, 벽은 로봇을 왼쪽으로 민다.
즉 벽이 로봇에 가하는 힘은 침투 방향의 반대이고, 로봇이 벽에 가하는 힘은 그 반대 방향이다.

임피던스 제어는 이 문제를 다른 방향에서 바라본다.
로봇이 특정 위치를 무조건 따라가게 만드는 대신, 로봇의 끝단이 마치 가상의 질량, 스프링, 감쇠기로 연결된 것처럼 행동하게 만든다.
설계자는 여기서 위치 오차를 힘으로 바꾸는 강성, 속도를 힘으로 바꾸는 감쇠, 그리고 필요한 경우 가속도에 대응하는 관성을 선택한다.
따라서 로봇은 목표점에 가까워지려는 성질을 가지되, 외부 환경과 접촉할 때는 적절히 밀리고 되돌아오는 물리적 반응을 보인다.
이 튜토리얼에서는 먼저 병진 질량-스프링-댐퍼 목표 동역학으로 직관을 잡고, 토크 제어, 동역학 보상, 센싱, 안정성, 액추에이터 제한 같은 실제 구현 조건은 뒤에서 구분한다.
물리적 상호작용이 필요한 로봇에서는 접촉력을 안전하고 능숙하게 활용해야 하며, 임피던스 제어는 구조화되지 않은 환경에서 큰 충격력을 피하기 위한 대표적인 접근으로 다루어져 왔다 \cite{abudakka2020variableimpedance}.

이 글은 새 제어 알고리즘을 내놓기보다, 전기, 진동, 로봇 제어에서 반복되는 임피던스의 감각을 입문자가 따라올 수 있게 풀어 쓰는 리뷰/튜토리얼이다.
성능 순위를 매기기보다 개념이 어떻게 이어지는지 보여주고, MuJoCo 실험실은 그 연결을 로그와 그래프로 다시 확인하는 작은 실험 지도 역할을 한다.
고급 제어 이론을 이미 알고 있다고 가정하지 않고, 수식도 엄밀한 증명보다 왜 필요한지와 플롯에서 어떤 신호로 확인되는지를 중심으로 읽는다.
따라서 MuJoCo 실습은 실제 하드웨어 성능 보증이나 정밀 디지털 트윈 검증이 아니라, 파라미터 변화가 응답의 어떤 특징으로 나타나는지 관찰하기 위한 개념 확인용 실습으로 보면 된다.
본문은 가능한 한 복잡한 수식을 앞세우기보다, 수식이 필요한 맥락과 물리적 의미를 먼저 설명한다.
왜 그런 수식이 필요해졌는지, 그 수식이 어떤 물리적 반응을 담고 있는지, 로봇 제어에서는 그 반응을 어떻게 설계 변수로 바꾸는지 순서대로 따라간다.
읽는 동안 계속 떠올릴 질문은 하나로 충분하다.
외부에서 밀리거나 흔들릴 때, 시스템은 얼마나 버티고, 얼마나 늦게 반응하며, 에너지를 어디에 저장하거나 잃는가?

이 질문은 전기회로에서는 전압과 전류의 관계로 나타나고, 기계진동에서는 힘과 운동의 관계로 나타나며, 로봇 제어에서는 접촉력과 끝단 움직임의 관계로 나타난다.
겉보기에는 서로 다른 분야처럼 보이지만, 세 경우 모두 에너지가 어디에 저장되고 어디에서 사라지는지를 추적하면 같은 구조가 드러난다.
처음에는 먼저 전압-전류, 힘-운동을 한 쌍으로 묶어 보면 된다.
다만 여기서 원인과 결과는 자연법칙의 한쪽 방향이 아니라, 우리가 어떤 실험을 하고 어떤 비율을 보려는지에 따라 정한 읽는 방향이다.
전기회로에서는 전압을 걸고 전류를 볼 수도 있지만, 전류원을 걸고 필요한 전압을 볼 수도 있다.
기계에서도 힘을 가해 속도나 변위를 볼 수도 있고, 정해진 운동을 만들려면 얼마나 큰 힘이 필요한지를 물을 수도 있다.
로봇 접촉에서는 외부 힘과 손끝 움직임이 서로 영향을 주고받는다.
임피던스 관점은 주로 운동에서 필요한 힘을 묻는 관점에 가깝고, 어드미턴스 관점은 힘이 주어졌을 때 생기는 운동을 묻는 관점에 가깝다.
두 관점 모두 같은 에너지 저장/소산 요소를 서로 다른 방향에서 읽는 방법이라고 두면 뒤의 제어 장이 더 선명해진다.

각 장을 읽을 때는 이 질문을 세 가지 작은 확인으로 나누어 보면 좋다.
첫째, 이 장에서 입력 또는 원인으로 보는 물리량은 무엇인가.
둘째, 출력 또는 반응으로 보는 물리량은 무엇인가.
셋째, 에너지를 저장하는 요소와 소산하는 요소는 각각 무엇인가.
이 세 가지가 보이면 전기회로, 진동계, 로봇 제어가 서로 다른 예시가 아니라 같은 구조가 분야마다 다른 모습으로 나타난다는 점이 보인다.

이 글은 주로 운동 상태를 보고 필요한 힘을 설계하는 임피던스 관점을 따라간다.
반대로 측정된 힘을 보고 로봇의 운동 명령을 만드는 관점은 어드미턴스 제어라고 부르며, 뒤의 임피던스 제어 장에서 다시 구분한다.
처음 읽을 때는 같은 힘-운동 관계를 어느 방향에서 읽고 구현하는지만 구분해도 충분하다.

\subsection{읽는 순서}

이 글의 흐름은 하나의 질문을 점점 구체화한다.
임피던스가 왜 필요했는지에서 출발해, 같은 힘-운동 관계를 회로, 1자유도 진동, 로봇 손끝으로 옮겨 읽는다.
처음부터 모든 용어를 완벽히 외울 필요는 없다.
각 장은 다음 장을 이해하기 위한 직관을 조금씩 마련한다.
각 장은 배경지식을 넓게 펼치기보다, 독자가 다음 장면을 조금 더 예측할 수 있게 설계되어 있다.
장을 마칠 때는 아래 질문에 스스로 답해 보면 된다.

\begin{enumerate}
  \item 먼저 임피던스가 왜 저항 하나로는 부족했던 문제에서 중요해졌는지 살펴본다. 이 튜토리얼에서는 긴 전보선과 전화선에서 신호가 뭉개지는 역사적 문제를 중요한 입구로 삼는다. 이 장면은 임피던스가 단순한 공식이 아니라 실제 공학 문제에서 자란 언어임을 보여준다. 독자는 저항, 리액턴스, 임피던스가 왜 구분되는지 말할 수 있어야 한다.
  \item 다음으로 선형시불변 시스템을 다룬다. 여기서 선형은 입력을 두 배로 키우면 출력도 두 배로 커지는 성질이고, 시불변은 같은 입력을 오늘 넣든 나중에 넣든 시스템의 규칙이 바뀌지 않는다는 뜻이다. 이는 전기회로, 진동계, 로봇 끝단을 같은 수학 언어로 비교하기 위한 최소한의 문법이다. 여기까지 오면 입력, 출력, 상태, 전달함수를 구분할 수 있어야 한다.
  \item 그다음 RLC 회로를 통해 저항, 인덕터, 커패시터가 각각 에너지를 어떻게 소산하고 저장하는지 본다. 이 장에서는 인덕터와 커패시터가 왜 전류와 전압의 변화를 방해하는지 설명할 수 있어야 한다.
  \item 이어서 질량-스프링-댐퍼 시스템으로 넘어가 같은 구조가 기계 시스템에서는 어떻게 나타나는지 확인한다. 여기까지 오면 질량, 스프링, 댐퍼가 고유진동수, 감쇠비, 오버슈트, 정착시간에 어떤 영향을 주는지 예측할 수 있어야 한다.
  \item 그다음 다관절 로봇을 읽기 위한 최소 문법을 준비한다. 관절공간과 작업공간, 정기구학과 역기구학, 자코비안, 힘--관절 effort 변환, 경로와 궤적을 알아야 1차원 스프링-댐퍼 직관이 실제 로봇 손끝으로 어떻게 옮겨지는지 볼 수 있다.
\vmark{scope-intro-force-effort-transform}%
  \item 이후 이 직관을 로봇 임피던스 제어로 옮긴다. 이때 평형점, 강성, 감쇠계수와 감쇠비는 설계 변수이고, 정상상태응답과 과도응답은 그 결과를 읽는 기준이 된다. 독자는 강성을 바꾸면 접촉력과 위치 오차가 어떻게 달라지고, 감쇠비를 바꾸면 진동과 정착시간이 어떻게 달라지는지 말할 수 있어야 한다.
  \item 끝으로 MuJoCo 기반 Lab01--Lab04에서 앞의 개념들이 로그, 그래프, 뷰어 화면으로 어떻게 나타나는지 정리한다. Lab04는 Franka Panda 모델의 위치 액추에이터 기반 가상 벽/목표점 후퇴 실험으로 읽고, 실제 접촉력 추출이나 하드웨어 안정성 검증으로 읽지 않는다. 마지막에는 정밀 로봇 검증이 아니라, 이론에서 예측한 응답 모양을 재현 가능한 설정과 그래프로 확인하는 방법을 익힌다.
\end{enumerate}

학습 산출물은 거창할 필요가 없다.
각 장을 읽은 뒤에는 핵심 개념을 한 문장으로 설명하고, 간단한 숫자 예 하나를 손으로 계산해 보고, 실험실 플롯에서 확인할 신호 하나를 고르면 충분하다.
예를 들어 저항과 임피던스의 차이를 말로 설명하고, $1~\mathrm{cm}$ 위치 오차가 어느 정도 힘이 되는지 계산하고, Lab 플롯에서 위치 오차와 제어 입력을 함께 확인하는 식이다.
뒤의 실험실 장에서는 이 흐름을 조금 더 구체적으로, \texttt{config.yaml}에서 하나의 값을 바꾸고 실행한 뒤 \texttt{log.csv}, \texttt{summary.json}, 필요한 경우 \texttt{--plot}으로 저장한 \texttt{plots/*.png}, \texttt{worksheet.md}를 차례로 보며 확인한다.
바꾸기 전에는 응답이 어떻게 달라질지 한 문장으로 예상하고, 실행 뒤에는 그 예상이 그래프와 로그에서 어디까지 맞았는지만 적어도 충분하다.

실험실 장을 읽을 때는 각 Lab을 별개의 데모로 보지 말고, 같은 질문이 커지는 순서로 보면 된다.
Lab01은 질량-스프링-댐퍼에서 강성, 감쇠, 관성이 시간응답을 어떻게 바꾸는지 확인한다.
Lab02는 목표 추종과 제어 입력 제한이 응답 지표에 어떤 변화로 나타나는지 본다.
다관절 로봇 기초 장은 이 직관을 관절공간, 작업공간, 자코비안, 경로와 궤적의 언어로 옮긴다.
Lab03은 같은 직관이 2자유도 로봇의 관절공간, 작업공간, 자코비안, 특이점 문제로 어떻게 확장되는지 보여준다.
Lab04는 7자유도 매니퓰레이터 모델에서 목표점, 계산된 가상 벽 신호, 목표점 후퇴, 강성-감쇠 선택을 로그와 플롯으로 짚는다.

이 순서를 지나면 독자는 단순히 용어를 외우는 데서 멈추지 않는다.
긴 전선에서 왜 신호가 뭉개졌는지, RLC 회로와 질량-스프링-댐퍼가 왜 같은 2차 구조를 가지는지, 강성과 감쇠비를 바꾸면 로봇 접촉 응답의 어느 부분이 달라지는지 예측할 수 있어야 한다.
끝의 실험실 장은 그 예측을 로그와 그래프에서 확인하는 방법을 보여준다.

그림~\ref{fig:impedance-roadmap}은 이 글의 전체 읽는 길을 한 장으로 요약한다.
이 그림은 임피던스의 모든 역사를 빠짐없이 정리한 연대기가 아니라, 입문자가 따라가기 쉬운 학습 흐름도이다.
처음에는 전기 통신선의 신호 왜곡 문제에서 출발하지만, 핵심 질문은 계속 같다.
어떤 요소가 에너지를 저장하고, 어떤 요소가 에너지를 소산하며, 그 결과 입력과 출력 사이의 관계가 어떻게 달라지는가를 보는 것이다.
그림의 화살표는 사건을 시간순으로 외우라는 표시가 아니라, 학습에 필요한 도구가 하나씩 넘어가는 길이다.
긴 전보선과 전화선은 저항만으로 시간에 따라 변하는 신호의 퍼짐, 지연, 위상 차이를 설명하기 어렵다는 문제를 보여준다.
복소 임피던스와 LTI 시스템은 그 문제를 크기, 위상, 입력-출력 관계로 정리하는 언어를 준다.
RLC 회로는 소산과 저장을 구체적인 전기 소자로 보여주고, 질량-스프링-댐퍼는 같은 저장과 소산 구조를 힘, 속도, 위치의 눈에 보이는 운동으로 옮긴다.
다관절 로봇 기초는 이 구조가 실제 관절, 손끝, 자코비안, 궤적을 거쳐 구현되어야 함을 보여준다.
로봇 임피던스 제어는 이 구조를 로봇 끝단에 가상으로 부여해 접촉 중 힘과 운동의 관계를 설계하고, MuJoCo 실험실은 그 선택이 로그와 플롯에서 어떤 모양으로 나타나는지 확인하는 마지막 단계이다.
다만 뒤의 Franka Panda 가상 벽 실험은 토크 수준 작업공간 임피던스 검증이 아니라, 위치 액추에이터와 결정론적 목표점 후퇴를 이용해 접촉력 직관을 보여주는 교육용 대리 실험으로 읽어야 한다.

\begin{figure}[!htbp]
  \centering
  \resizebox{\linewidth}{!}{%
  \begin{tikzpicture}[
      stage/.style={
        draw=blue!55!black,
        rounded corners=2pt,
        fill=blue!4,
        align=center,
        text width=3.0cm,
        minimum height=1.35cm,
        inner sep=5pt,
        font=\small
      },
      note/.style={
        draw=black!35,
        rounded corners=2pt,
        fill=black!3,
        align=center,
        text width=2.95cm,
        inner sep=4pt,
        font=\scriptsize
      },
      arr/.style={-{Latex[length=2.2mm]}, thick, draw=black!55}
    ]
    \node[stage] (line) {긴 전보선ㆍ전화선\\신호 왜곡};
    \node[stage, right=6mm of line] (complex) {복소 임피던스\\$V=IZ$};
    \node[stage, right=6mm of complex] (rlc) {RLC 회로\\저장과 소산};
    \node[stage, right=6mm of rlc] (msd) {질량-스프링-댐퍼\\기계 임피던스};
    \node[stage, right=6mm of msd] (robot) {로봇 끝단\\가상 임피던스};
    \node[stage, right=6mm of robot] (lab) {MuJoCo Lab\\로그와 plot};

    \draw[arr] (line) -- (complex);
    \draw[arr] (complex) -- (rlc);
    \draw[arr] (rlc) -- (msd);
    \draw[arr] (msd) -- (robot);
    \draw[arr] (robot) -- (lab);

    \node[note, below=5mm of line] {왜 저항만으로\\설명이 부족한가?};
    \node[note, below=5mm of complex] {크기와 위상을\\함께 본다};
    \node[note, below=5mm of rlc] {$R$: 소산\\$L,C$: 저장};
    \node[note, below=5mm of msd] {$b$: 소산\\$m,k$: 저장};
    \node[note, below=5mm of robot] {힘과 운동의 관계를\\어떻게 설계하는가?};
    \node[note, below=5mm of lab] {파라미터 변화가\\응답을 어떻게 바꾸는가?};
  \end{tikzpicture}%
  }
  \caption{이 튜토리얼이 따라가는 임피던스 개념의 이동 경로. 이 그림은 임피던스의 완전한 역사 연대기가 아니라, 긴 통신선의 신호 왜곡이라는 중요한 역사적 동기에서 출발해 에너지 저장과 소산, 기계적 임피던스, 로봇 끝단의 접촉 반응 설계까지 이어지는 학습 경로를 요약한 것이다.}
  \label{fig:impedance-roadmap}
\end{figure}


이 흐름에서 중요한 것은 분야 이름이 아니라 역할이다.
저항과 댐퍼는 에너지를 소산하고, 인덕터와 질량은 변화에 대한 관성을 만들며, 커패시터와 스프링은 위치 또는 상태에 따른 에너지를 저장한다.
임피던스 제어는 이 역할들을 로봇 끝단에 가상으로 부여하는 방법으로 이해할 수 있다.

\subsection{처음 읽을 때 필요한 용어 지도}

이 글에는 전기, 기계, 제어, 로봇 용어가 함께 등장한다.
처음부터 모두 외울 필요는 없다.
다음 용어만 붙잡고 있으면 뒤에서 나오는 수식이 어떤 물리량을 다루는지 따라가기 쉽다.
이 표는 암기표라기보다, 수식을 읽다가 기호가 헷갈릴 때 다시 확인하는 작은 참조표라고 생각하면 된다.
기호가 흔들릴 때는 부록~\ref{sec:notation-checklist}의 빠른 참조표로 돌아가면 된다.
특히 $q$, $C$, $V$, $E$처럼 장마다 뜻이 달라질 수 있는 기호는 아래첨자와 문맥을 먼저 확인한다.

\begin{table}[!htbp]
  \centering
  \caption{핵심 용어 지도}
  \begin{tabular}{p{0.22\linewidth}p{0.34\linewidth}p{0.32\linewidth}}
    \toprule
    용어 & 쉬운 의미 & 이 글에서의 역할 \\
    \midrule
    순응성/compliance & 힘을 받았을 때 얼마나 잘 밀리는가 & 정적 강성의 역수에 가까운 개념 \\
    강성/stiffness & 얼마나 안 밀리고 버티는가 & 위치 오차를 힘으로 바꿈 \\
    감쇠/damping & 움직임의 에너지를 얼마나 줄이는가 & 진동과 오버슈트를 줄임 \\
    관성/inertia & 속도 변화에 얼마나 둔감한가 & 가속과 충격 응답을 결정 \\
    작용 변수/effort & 시스템을 밀거나 누르는 원인 변수 & 전압, 힘, 접촉력을 같은 역할로 비교 \\
    흐름 변수/flow & 원인에 의해 흐르거나 움직이는 반응 변수 & 전류, 속도, 끝단 운동을 같은 역할로 비교 \\
    주파수 영역 & 빠른 변화와 느린 변화를 나누어 보는 관점 & 임피던스와 공진을 설명 \\
    전달함수 & 입력과 출력의 관계를 담은 함수 & 응답 모양을 예측 \\
    과도응답 & 목표나 외란 뒤에 흔들리고 정착하는 중간 과정 & 오버슈트, 진동, 정착시간을 읽는 구간 \\
    정상상태응답 & 시간이 충분히 지난 뒤 남는 응답 & 최종 위치 오차와 일정 접촉력의 관계를 읽는 구간 \\
    감쇠비 & 감쇠가 진동을 얼마나 줄이는지 보는 무차원 비율 & 감쇠계수 하나만으로 판단하지 않게 해줌 \\
    관절공간 & 로봇 관절각으로 보는 좌표 & 관절각으로 명령을 쓰는 방식 \\
    작업공간 & 로봇 손끝 위치로 보는 좌표 & 손끝 위치로 명령을 쓰는 방식 \\
    \bottomrule
  \end{tabular}
\end{table}

순응성은 특히 중요하다.
강성이 크다는 말은 같은 힘에 대해 적게 밀린다는 뜻이고, 순응성이 크다는 말은 같은 힘에 대해 많이 밀린다는 뜻이다.
스프링 하나만 보면 강성은 $k$, 순응성은 대략 $1/k$이다.
강성의 단위는 보통 $\mathrm{N/m}$이고, 순응성의 단위는 그 역수인 $\mathrm{m/N}$이다.
즉 순응성은 힘 $1~\mathrm{N}$을 가했을 때 얼마나 많이 밀리는지를 보는 정적 감각에 가깝다.
동적인 힘-운동 관계의 역방향은 뒤에서 어드미턴스라는 이름으로 따로 구분한다.
부드러운 로봇이라고 말할 때는 강성이 낮거나 순응성이 높다고 표현할 수 있다.

이 글은 두 층으로 구성된다.
첫 번째 층은 공통 언어이다.
임피던스, LTI 시스템, RLC 회로, 질량-스프링-댐퍼를 차례로 보면서 강성, 감쇠, 관성, 에너지 저장과 소산을 같은 문법으로 읽는다.
두 번째 층은 실험실이다.
수식으로 이해한 개념이 실제 로그와 그래프에서 어떤 모양으로 나타나는지 확인한다.

이때 시뮬레이터는 이론을 대신하는 단순한 데모가 아니라, 이론의 감각을 반복해서 확인하는 실험 도구이다.
MuJoCo는 접촉과 관절 동역학을 포함하는 물리 엔진으로 널리 사용되어 왔으며 \cite{todorov2012mujoco}, DeepMind Control Suite와 dm\_control 같은 소프트웨어 계층에서도 동일 모델과 설정을 저장했을 때 반복 가능한 로컬 제어 실험을 구성하는 데 활용되었다 \cite{tassa2018deepmindcontrolsuite,tassa2020dmcontrol}.
로봇 조작 분야에서는 robosuite, MuJoCo Playground, Franka 기반 MuJoCo 환경처럼 다양한 시뮬레이션 환경이 제안되어 왔다 \cite{zhu2020robosuite,zakka2025mujocoplayground,xu2024frankamujocoenvs}.
이 논문에서 다루는 실험실은 이러한 대규모 벤치마크와 달리, 새 제어 알고리즘의 우열을 증명하기보다 한 학생이 로컬 컴퓨터에서 실행하고 그래프로 확인할 수 있는 작고 읽기 쉬운 개념 확인용 실습 환경에 초점을 둔다.
특히 Franka Panda 가상 벽 예시는 실제 접촉력 측정, 하드웨어 안전성, 토크 수준 임피던스 제어 검증을 목표로 하지 않는다.
위치 액추에이터, 계산된 가상 벽 신호, 목표점 후퇴를 통해 강성-감쇠 선택의 효과를 관찰하는 교육용 실험으로 둔다.
순응성, 어드미턴스, 강성, 임피던스처럼 서로 가까운 용어는 중간에 혼동될 수 있으므로, 뒤의 논의 장에서 자주 생기는 오해를 다시 표로 정리한다.
실험실 이야기는 마지막 장에서 다시 자세히 다룬다.
그 전에 먼저 임피던스를 읽는 공통 언어를 만들어 두자.

